%% Figure captions, supplied as a separate file (Applied Soft Computing, figures section).
%% Generated by scripts/build_latex.py from manuscript/sections/12_figure_captions.md.

\noindent\textbf{Figure 1.} \textbf{The arm lattice.} Each arrow is a single-factor contrast and the label names the one thing it varies; the other configured inputs are held fixed by a run-time fingerprint over the solver configuration, realizer, weight grid, surrogate identities and reference. The three primary contrasts are marked.\par\medskip
\noindent\textbf{Figure 2.} \textbf{The geometry contrast, per replication.} Paired differences in CORE-relative hypervolume ratio, NBI-S minus WS-S, sorted within each dataset, $R = 30$, CORE reference. Blue favours NBI-S. The dashed line is the median. Note that the vertical scales differ by dataset.\par\medskip
\noindent\textbf{Figure 3.} \textbf{Anchor provenance and the associated geometry.} Left: the full NBI-S $\rightarrow$ NBI-R difference beside the anchor-injection control, which inserts NBI-R's empirical anchors into NBI-S's set and changes nothing else. Right (post hoc): the CHIM extent under empirical anchors relative to surrogate-derived anchors, as a per-dataset median. The juxtaposition is a post hoc association, not a demonstrated cause: clear on Adult and Bank Marketing, weak on Spambase and absent on MAGIC.\par\medskip
\noindent\textbf{Figure 4.} \textbf{Reference sensitivity.} Median paired difference under the CORE reference (primary, horizontal) against the AUGMENTED reference (mandatory sensitivity, vertical), for all twelve dataset-by-contrast cells. Red marks the single cell whose significance differs between references; no cell changes direction.\par\medskip
\noindent\textbf{Figure 5.} \textbf{Dispersion across the panel.} The same paired differences as Figure 2, shown as distributions. The boundary dataset's spread is roughly thirty times MAGIC's, which is why its contrast does not resolve at $R = 30$ despite a numerically larger point estimate.\par\medskip
\noindent\textbf{Figure 6.} \textbf{Surrogate-gate regime and baselines.} Left: per-dataset pass rates of the frozen external reliability criterion over $R = 30$; the gate is diagnostic and filters nothing. Right: median CORE-relative hypervolume ratio by method. The dashed rule marks the CORE reference itself and is \emph{not} an upper bound: a ratio above it means the candidate set improved on the finite reference.\par\medskip
